\conclusion

Выпускная квалификационная работа посвящена теме управления инфраструктурой. В рамках исследования рассмотрена проблема поддержания соответствия между архитектурным описанием инфраструктуры и ее фактическим состоянием, а также разработана платформа, реализующая базовый цикл управления.

Актуальность выбранной темы связана с тем, что современная инфраструктура редко остается статичной. Управляемые хосты, контейнерные workload, средства наблюдаемости и вспомогательные сервисы изменяются в результате плановых релизов, ручных операций, сбоев, неполных процедур развертывания и изменения требований эксплуатации. В таких условиях архитектурная документация быстро теряет связь с реальностью, если она используется только как описание для человека. Подход Architecture-as-Code позволяет рассматривать архитектурную модель как формализованный артефакт, пригодный для машинной обработки, однако сам по себе он не решает задачу контроля фактического состояния. Поэтому в работе была исследована возможность использовать архитектурное описание как источник desired state для автоматизированной платформы управления инфраструктурой.

В ходе исследования были рассмотрены теоретические основы управления инфраструктурой в распределенных системах, понятия desired state, actual state, drift и reconcile, а также близкие подходы Infrastructure-as-Code и Architecture-as-Code. Был выполнен аналитический обзор существующих решений, который показал, что существующие инструменты решают важные части задачи, но не закрывают полностью выбранный сценарий. Это обосновало необходимость проектирования собственной платформы.

В результате проектирования были сформированы функциональные и нефункциональные требования к платформе, определены внешние участники и выполнена контейнерная декомпозиция. Архитектура решения построена вокруг четырех основных сервисов, компоненты имеют разделенные зоны ответственности и взаимодействуют через API-контракты, что упрощает развитие отдельных частей системы. ParserSvc отделяет обработку архитектурного источника от логики сравнения состояний. InventorySvc изолирует получение сведений о фактической инфраструктуре. DriftDetectorSvc концентрирует правила обнаружения расхождений и принятия решения о необходимости reconcile. ReconcilerSvc отделяет управляющее решение от непосредственного выполнения восстановительных действий. Такая структура позволяет рассматривать платформу не как набор отдельных скриптов, а как основу для дальнейшего control plane управления инфраструктурными workload.

В ходе реализации были поддержаны два базовых типа workload: \texttt{node\_exporter} и \texttt{cadvisor}. Выбор этих компонентов связан с их ролью в контуре наблюдаемости: отсутствие подобных workload может приводить к потере метрик и снижению диагностируемости инфраструктуры. На этом сценарии была проверена ключевая идея работы: архитектурное описание задает ожидаемый состав инфраструктурных компонентов, inventory-слой показывает их фактическое наличие, детекторы обнаруживают расхождение, а reconcile-контур инициирует восстановление отсутствующего workload.

Практическая значимость работы состоит в снижении объема ручных операций при контроле состояния инфраструктуры и в повышении наблюдаемости расхождений между архитектурным описанием и фактической средой. Для SRE- и DevOps-команд предложенный подход полезен тем, что связывает архитектурную модель, inventory, drift detection и reconcile в единый процесс. Это создает основу для более системного управления инфраструктурой, в котором расхождения выявляются не только в результате инцидентов или ручных проверок, но и в рамках регулярного программного цикла.

При этом разработанное решение имеет границы MVP. Текущая версия ориентирована на ограниченный набор workload, использует cAdvisor как основной источник фактических данных о контейнерных компонентах, хранит часть состояния в памяти и не реализует полноценную персистентность истории detection cycle и reconcile-задач. Также отсутствуют зрелые механизмы аутентификации сервисов, авторизации действий, политики допуска операторов. Эти ограничения не отменяют полученного результата, но определяют область применимости текущей реализации и показывают, какие доработки необходимы для промышленного использования.

Основными направлениями дальнейшего развития являются интеграция с NetBox или аналогичными системами учета инфраструктуры, развитие метрик и наблюдаемости самой платформы, добавление новых детекторов и операторов, а также усиление безопасности. Отдельным перспективным направлением является развитие Architecture-as-Code-контура: валидация архитектурных манифестов, поддержка нескольких источников описания, версионирование desired state и анализ изменений между версиями модели.

Таким образом, цель выпускной квалификационной работы достигнута. Разработана и обоснована автоматизированная платформа управления инфраструктурой с поддержкой подхода Architecture-as-Code, обеспечивающая формирование desired state из архитектурного описания, получение actual state, обнаружение drift и запуск reconcile для поддерживаемых инфраструктурных workload. Полученный результат подтверждает реализуемость предложенного подхода и создает основу для дальнейшего развития платформы в сторону более зрелого внутреннего control plane для управления инфраструктурой.
